\documentclass[conference]{IEEEtran}
\IEEEoverridecommandlockouts
\usepackage{cite}
\usepackage{amsmath,amssymb,amsfonts}
\usepackage{graphicx}
\usepackage{textcomp}
\usepackage{xcolor}
\usepackage{booktabs}
\usepackage{microtype}
\usepackage[hidelinks]{hyperref}

\def\BibTeX{{\rm B\kern-.05em{\sc i\kern-.025em b}\kern-.08em
    T\kern-.1667em\lower.7ex\hbox{E}\kern-.125emX}}

\newcommand{\system}{Socrates AI}

\begin{document}

\title{\system: A Bounded-Latency Virtual Mentor for Real-Time Personalized Programming Feedback in CS1}

\author{
    \IEEEauthorblockN{Alimzhan Koshikov}
    \IEEEauthorblockA{
        \textit{Master Degree Student, School of Software Engineering}\\
        Astana IT University (AITU)\\
        Astana, Kazakhstan\\
        225148@astanait.edu.kz\\
        \href{https://orcid.org/0009-0003-3234-6487}{0009-0003-3234-6487}
    }
    \and
    \IEEEauthorblockN{Tutkyshbayeva Shyryn}
    \IEEEauthorblockA{
        \textit{PhD, School of Software Engineering}\\
        Astana IT University (AITU)\\
        Astana, Kazakhstan\\
        sh.tutkyshbayeva@astanait.edu.kz\\
        \href{https://orcid.org/0000-0002-1419-2571}{0000-0002-1419-2571}
    }
}

\maketitle

\begin{abstract}
Introductory programming students often receive feedback that is delayed, generic, or limited to correctness checks. As a result, repeated syntax and logic errors are not always translated into pedagogically useful guidance during active problem solving. This paper presents \system, a virtual mentor architecture for CS1 that delivers bounded-latency personalized feedback by combining continuous code analysis, a lightweight student state, a pedagogical policy selector, and constrained feedback realization. The technical contribution is not a new foundation model; rather, it is a systems design that explicitly separates analyzer-derived features from policy selection and language realization, enabling both a rule-based baseline and an optional learned policy over the same feature space. The implemented prototype records interaction logs for later policy learning while keeping feedback generation under explicit pedagogical guardrails. A pilot event-level evaluation on 60 logged interactions showed a mean end-to-end latency of \textbf{741\,ms} and \textbf{809\,ms} at P95, with an event-level success rate of \textbf{66.7\%} and a median time-to-fix of \textbf{2.45\,s}. These results support the feasibility of a real-time virtual mentor pipeline for CS1, while larger classroom-scale and longitudinal evaluations remain future work.
\end{abstract}

\begin{IEEEkeywords}
virtual mentor, CS1, introductory programming, programming education, real-time feedback, student modeling, pedagogical policy, bounded latency
\end{IEEEkeywords}

\section{Introduction}

Introductory programming (CS1) remains difficult for many first-year students because they must learn syntax, program logic, and problem decomposition at the same time. In large classes, beginners often move through repeated cycles of compiler errors, failing tests, and local misconceptions while instructor attention is limited. Timely formative feedback is therefore important, yet individualized support does not scale well in typical teaching settings \cite{DesirableCharacteristicsforAITeachingAssistantsinProgrammingEducation,FeasibilityStudyofAugmentingTeachingAssistantswithAI}.

Automated assessment partially addresses this scalability problem through unit tests, static checks, and compiler diagnostics. However, such feedback is often binary and rarely adapts to the learner's current difficulty, repeated mistakes, or recent feedback history. Research in computing education shows that useful AI assistance depends not only on correctness, but also on timing, pedagogical framing, and whether the interaction preserves student autonomy instead of encouraging answer-seeking behavior \cite{SystematicLiteratureReviewonLLM,PatternsofLLMPoweredProgrammingAssistant,DesirableCharacteristicsforAITeachingAssistantsinProgrammingEducation}. Recent deployments such as CS50's AI support further show that near real-time help is technically feasible, but maintaining reliability and pedagogical alignment requires explicit evaluation and continuous refinement \cite{ImprovingAIinCS50}.

Existing work still leaves a practical gap. Intelligent tutoring systems offer learner adaptation but often rely on heavier domain engineering or intervene after discrete checkpoints rather than during active code writing \cite{ProgTutor,Francisco2022}. LLM-based assistants can generate fluent hints and explanations, but many operate as general-purpose conversational tools without explicit learner-state modeling or fine-grained control over the timing and disclosure level of feedback \cite{CodeAid,IRIS,TheProgrammersAssistant}. For CS1, this matters: classroom evidence suggests that unconstrained or overly direct assistance can reduce conceptual engagement and lead to over-reliance \cite{PatternsofLLMPoweredProgrammingAssistant,FeasibilityStudyofAugmentingTeachingAssistantswithAI}.

This paper presents \system, the prototype implementation used in the author's master's thesis on virtual mentor models for programming education. The system is designed as a bounded real-time pipeline with four explicit layers: (i) code analysis, (ii) lightweight student-state construction, (iii) pedagogical policy selection, and (iv) constrained feedback realization. The key novelty is a system-level integration that treats policy selection as a separate component from both analysis and text generation. This makes the architecture interpretable, allows a hand-written rule baseline to coexist with a learned selector, and keeps the LLM---when used---limited to wording rather than pedagogical decision-making.

The contributions of this paper are therefore:
\begin{itemize}
\item a modular virtual mentor architecture for CS1 that separates analyzer signals, student-state features, pedagogical policy selection, and constrained feedback realization;
\item an implementation of this architecture in \system, including a rule-based baseline policy, an optional learned policy interface, interaction logging for future supervised policy learning, and explicit guardrails against direct-solution disclosure;
\item a pilot event-level evaluation of the prototype, reporting bounded-latency characteristics and action-level effectiveness trends while clearly delimiting the current evidence to prototype-scale results rather than classroom-scale learning claims.
\end{itemize}

\section{Related Work and Positioning}

Research on AI support for programming education can be grouped into three relevant directions: intelligent tutoring systems, automated feedback systems, and LLM-based assistants.

\subsection*{A. Intelligent Tutoring Systems}

Intelligent tutoring systems (ITS) adapt support based on learner state, task progression, or instructional sequencing. Recent examples such as ProgTutor emphasize augmentation and personalized sequencing rather than full instructional automation \cite{ProgTutor}. This line of work provides important ideas for learner modeling, but many ITS solutions either require substantial domain engineering or focus on step-level progression rather than low-latency feedback during free-form code writing \cite{Francisco2022,Chrysafiadi2022}.

\subsection*{B. Automated Feedback Systems}

Automated feedback systems scale well and are widely used in programming education. Typical implementations rely on compiler diagnostics, tests, program analysis, or repair generation \cite{AutomatedAssessmentToolsForGrading2022,AutomatedGradingandFeedbackToolsforProgrammingEducationSystematicReview,NavigatingAutomatedFeedback}. Their main strength is operational efficiency; their main limitation is that they often focus on correctness or repair without adapting the intervention to learner history or pedagogical intent.

\subsection*{C. LLM-Based Assistants}

LLM-based programming assistants offer flexible explanation and hint generation in natural language \cite{SystematicLiteratureReviewonLLM,pankiewicz2023largelanguagemodelsgpt}. Systems such as CodeAid and Iris demonstrate that this can be pedagogically useful when guardrails are added \cite{CodeAid,IRIS}. However, prior classroom studies also show the risk of over-reliance, shallow help-seeking, and explanation quality problems when the assistant is treated as a general answer source rather than a pedagogically controlled tool \cite{PatternsofLLMPoweredProgrammingAssistant,FeasibilityStudyofAugmentingTeachingAssistantswithAI}.

\subsection*{D. Positioning of the Proposed System}

The proposed work is positioned as a \emph{systems and prototype paper}. Its scientific contribution is primarily architectural and operational: it combines real-time analyzer signals, lightweight learner-state tracking, policy-level pedagogical control, and bounded-latency realization in one deployable pipeline. It does \emph{not} claim a new foundation model or a classroom-scale proof of learning gains. Instead, it contributes a concrete virtual mentor implementation that makes policy selection explicit and measurable.

Table~\ref{tab:positioning} clarifies how the proposed system differs from representative prior approaches.

\begin{table*}[t]
\centering
\caption{Positioning of the proposed system relative to representative prior work}
\label{tab:positioning}
\renewcommand{\arraystretch}{1.15}
\begin{tabular}{p{2.7cm} p{2.3cm} p{2.2cm} p{2.8cm} p{2.7cm} p{2.6cm}}
\hline
\textbf{System} & \textbf{Primary Mode} & \textbf{In-Process Feedback} & \textbf{Explicit Student State} & \textbf{Pedagogical Action Control} & \textbf{Latency-Centered Prototype Reporting}\\
\hline
AIF 3.0 \cite{AdaptiveImmediateFeedbackforBlockBasedProgramming} & Adaptive feedback for block-based tasks & Yes & Partial & Partial & Yes\\
ProgTutor \cite{ProgTutor} & ITS / adaptive tutoring & Partial & Yes & Yes & Limited\\
Iris \cite{IRIS} & LLM-based virtual tutor & Yes & Limited & Partial & Limited\\
CodeAid \cite{CodeAid} & LLM-based classroom assistant & Yes & Limited & Yes & Limited\\
\system{} (this paper) & Real-time virtual mentor pipeline & Yes & Yes & Yes & Yes\\
\hline
\end{tabular}
\end{table*}

\section{\system{} Architecture}

\subsection*{A. Design Objective}

\system{} is designed to answer one systems question: \emph{how can a CS1 mentor provide personalized and pedagogically constrained feedback in real time without relying on an unconstrained LLM to decide what to say?} To answer this, the implementation separates the decision pipeline into four layers so that each part is observable, replaceable, and measurable.

\begin{figure}[t]
\centering
\fbox{\parbox{0.95\linewidth}{
\centering
\textbf{Student code edits / submission}\\[3pt]
$\downarrow$\\[3pt]
\textbf{Analyzer}\\
error type, severity, compile status, test results, suspicious region, code lines\\[3pt]
$\downarrow$\\[3pt]
\textbf{Student-state builder}\\
error recurrence, attempt number, last feedback action, last feedback success, session feedback count\\[3pt]
$\downarrow$\\[3pt]
\textbf{Policy selector}\\
rule-based baseline or learned classifier over the same feature interface\\[3pt]
$\downarrow$\\[3pt]
\textbf{Guardrails + feedback realization}\\
template or LLM phrasing under no-solution-disclosure constraints
}}
\caption{Operational pipeline of \system. The pedagogical decision is separated from text realization.}
\label{fig:architecture}
\end{figure}

\subsection*{B. Operational Pipeline}

The analyzer produces structured signals from the current code state rather than natural-language descriptions. In the prototype, these signals include error type, severity, compile success, passed/failed tests, suspicious code region, analysis time, and code length. A separate student-state layer augments these analyzer outputs with short-horizon interaction context such as recurrence of the same error, total errors seen in the current session, attempt number, prior feedback action, whether the previous feedback helped, and the number of feedback events already issued.

This explicit feature interface is central to the architecture. It allows the pedagogical policy to operate on a compact and interpretable state representation instead of raw code or free-form dialogue. It also keeps the system extensible: the same feature vector can feed either a hand-written baseline policy or a learned classifier trained from interaction logs.

\subsection*{C. Formal Policy View}

Let $c_t$ denote the student code state after an edit at time $t$. The analyzer maps code to a structured signal vector,
\[
A: c_t \rightarrow z_t,
\]
where $z_t$ includes the current error type, severity, compile result, test outcomes, and localization hints.

A lightweight student state $s_t$ accumulates recent interaction context such as recurrence counts and prior feedback effects. The pedagogical policy then selects an action
\[
\pi: (z_t, s_t) \rightarrow a_t,
\]
where
\[
\mathcal{A}=\{\textsc{CodeHighlight},\ \textsc{ConceptualHint},\ \textsc{GuidingQuestion},\ \textsc{NoFeedback}\}.
\]

In the baseline configuration, $\pi$ is a hand-written decision function over the feature vector. In the extended configuration, the same interface supports a learned multi-class selector,
\[
\hat{a}=\arg\max_{a \in \mathcal{A}} P(a \mid x_t),
\]
where $x_t$ concatenates analyzer signals and student-state features. This separation is deliberate: machine learning is applied to \emph{policy selection}, while language generation remains only a realization step.

\subsection*{D. Baseline Policy and Learned Policy Extension}

The current prototype includes a rule-based baseline that maps recurring novice-error patterns to pedagogical actions. For example, syntax errors and first occurrences of off-by-one mistakes are mapped to code highlighting, while repeated condition-related mistakes are escalated to conceptual hints and no-progress situations trigger guiding questions. A learned policy interface is then defined over the same feature schema so that interaction logs can later be used for supervised action prediction. This design directly supports a scientifically defensible comparison: the rule engine serves as the baseline policy and the learned selector becomes an extension rather than a replacement of the whole system.

\subsection*{E. Realization Layer and Guardrails}

The realization layer converts a selected action into short learner-facing feedback. In \system{}, template-based feedback is always available, while LLM phrasing is optional. Crucially, the LLM does not choose the action; it only verbalizes a preselected pedagogical action under explicit prompt-level constraints such as ``do not reveal the full solution'' and ``keep the response concise.'' This design follows recommendations from prior work on educational AI assistants that stress the importance of preserving learner autonomy and avoiding direct solution disclosure \cite{DesirableCharacteristicsforAITeachingAssistantsinProgrammingEducation,CodeAid,AutomatingHumanTutorStyleProgrammingFeedback}.

\subsection*{F. Implementation in the Thesis Prototype}

The implemented thesis prototype, \system{}, realizes this pipeline as a web-based system. A Spring Boot backend orchestrates analyzer, student-state, policy, and feedback services; a frontend provides the coding workspace and displays feedback; and an optional Python microservice can host a learned policy model. Interaction events are stored in a structured log containing feature values, selected feedback actions, latency components, and post-feedback outcomes. This logging strategy is part of the technical contribution because it creates a reusable dataset for future policy learning and comparative evaluation.

\section{Pilot Evaluation}\label{sec:evaluation}

\subsection*{A. Evaluation Goal and Artifact}

The current evaluation is a \emph{pilot prototype evaluation}, not a classroom-scale effectiveness study. Its purpose is to answer two narrower questions:
\begin{itemize}
\item Can the architecture operate within interactive time bounds during use?
\item Do the logged feedback actions show useful event-level resolution trends?
\end{itemize}

The evaluation artifact is an event-level interaction log collected from prototype use on introductory programming tasks. The analyzed pilot trace contains \textbf{60} logged feedback events. Because the present artifact is an event stream rather than a balanced classroom trial dataset, the paper reports system-level and event-level measures only. It does not claim statistically validated learning gains, benchmark superiority, or longitudinal retention effects.

Each event record stores: detected error type, severity, compile status, test outcomes, recurrence counts, prior feedback context, selected pedagogical action, latency components, and post-feedback resolution indicators. This level of reporting directly addresses the reviewer concern that the dataset should be described in technical rather than purely pedagogical terms.

\subsection*{B. Metrics}

The prototype is evaluated with both operational and event-level feedback metrics.

The end-to-end response time is defined as
\[
T_{total} = T_{analysis} + T_{policy} + T_{feedback} + T_{overhead},
\]
where $T_{analysis}$ is analyzer time, $T_{policy}$ is pedagogical decision time, $T_{feedback}$ is realization time, and $T_{overhead}$ captures remaining platform latency.

The mean system latency is then
\[
\bar{T}=\frac{1}{N}\sum_{i=1}^{N} T_{total}^{(i)}.
\]

For event-level effectiveness, a feedback event is marked successful if the detected issue is resolved after that feedback event. We also report time-to-fix to estimate how quickly the learner moved from the flagged state to a resolved one.

\subsection*{C. Results}

Table~\ref{tab:results} summarizes the pilot trace. The main finding is that the system remained below one second at both mean and P95 latency, which is sufficient for interactive coding support.

\begin{table}[t]
\centering
\caption{Pilot results from event-level interaction logs}
\label{tab:results}
\renewcommand{\arraystretch}{1.15}
\begin{tabular}{lc}
\hline
\textbf{Metric} & \textbf{Value} \\
\hline
Events ($N$) & 60 \\
Mean latency $\bar{T}$ (ms) & 741 \\
P95 latency (ms) & 809 \\
Mean analysis/policy/feedback (ms) & 618 / 26 / 12 \\
Platform overhead (ms) & 85 \\
Event success rate (\%) & 66.7 \\
Median time-to-fix (ms) & 2450 \\
P95 time-to-fix (ms) & 3595 \\
\hline
\end{tabular}
\end{table}

Table~\ref{tab:action_success} shows action-level resolution rates in the pilot trace. In this prototype sample, code highlighting and conceptual hints were associated with stronger immediate resolution rates than guiding questions, suggesting that novice-error situations in the current tasks often benefited from slightly more explicit intervention.

\begin{table}[t]
\centering
\caption{Event success rate by feedback action in the pilot trace}
\label{tab:action_success}
\renewcommand{\arraystretch}{1.15}
\begin{tabular}{lcc}
\hline
\textbf{Action} & \textbf{N} & \textbf{Success rate} \\
\hline
CODE\_HIGHLIGHT & 23 & 1.00 \\
CONCEPTUAL\_HINT & 15 & 1.00 \\
GUIDING\_QUESTION & 19 & 0.25 \\
NO\_FEEDBACK & 3 & 0.00 \\
\hline
\end{tabular}
\end{table}

These results should be interpreted carefully. They demonstrate prototype feasibility and provide action-level operational insight, but they are not a substitute for a controlled benchmark comparison or student-level outcome study.

\section{Discussion}

\subsection*{A. What the Current Results Actually Support}

The pilot results support three claims.

First, the architecture is operationally viable: the full analyzer--policy--realization pipeline remained within interactive time bounds. Second, explicit policy separation is useful because it makes action-level behavior measurable rather than hiding pedagogical decisions inside free-form generation. Third, the event trace suggests that different feedback actions behave differently in practice, which justifies treating policy selection as a first-class research problem instead of assuming that one generic feedback style is sufficient.

These are meaningful technical contributions for a systems-oriented paper. They show that real-time personalized mentoring can be implemented as a measurable control pipeline rather than as an unconstrained chatbot wrapper.

\subsection*{B. Addressing the Reviewers' Concerns}

The review process raised four legitimate concerns: novelty positioning, technical depth, dataset clarity, and evaluation scope. The revised manuscript addresses them as follows.

Novelty is now stated explicitly as a \emph{systems architecture contribution}: analyzer signals, learner-state features, policy selection, and text realization are treated as separate layers. Technical depth is strengthened by formally describing the policy interface and by detailing how a rule baseline and a learned selector can operate over the same feature representation. Dataset clarity is improved by describing the evaluation artifact as a 60-event interaction log with explicit operational fields rather than as an underspecified pedagogical experiment. Finally, evaluation claims are deliberately narrowed to prototype feasibility and event-level trends, with larger classroom studies framed as future work instead of implied results.

\subsection*{C. Limitations}

The limitations remain important.
\begin{itemize}
\item The present evaluation is pilot-scale and event-level; it does not support inferential claims about long-term learning gains.
\item The prototype has not yet been benchmarked against a classroom-scale non-personalized baseline or an external tutoring system under controlled conditions.
\item The current learner model is intentionally lightweight. This helps with interpretability and latency, but it does not capture richer temporal learning dynamics.
\item The learned policy path is architecturally prepared and log-compatible, but this paper does not yet report a full baseline-vs-ML comparison.
\end{itemize}

These limitations do not invalidate the prototype contribution, but they do bound the scope of the conclusions.

\subsection*{D. Practical Implications}

From a practical perspective, the results suggest that a virtual mentor can be positioned between static automated checks and human tutoring. In this role, the system is not expected to replace instructors. Instead, it provides low-latency pedagogically constrained support for recurring novice difficulties and records interaction evidence that can later inform stronger personalization models. This aligns with prior arguments that AI assistants in programming education are most useful when they complement teaching workflows rather than operate as unrestricted solution generators \cite{CodeAid,ImprovingAIinCS50,AITutoringinSDEducation}.

\section{Conclusion}

This paper presented \system, a bounded-latency virtual mentor architecture for introductory programming. The implemented prototype separates code analysis, lightweight student-state construction, pedagogical policy selection, and constrained feedback realization. This architectural separation is the main technical contribution of the work because it creates an interpretable and extensible mentoring pipeline in which policy selection can be studied independently from language generation.

A pilot log-based evaluation on 60 feedback events showed mean end-to-end latency of \textbf{741\,ms}, P95 latency of \textbf{809\,ms}, an event-level success rate of \textbf{66.7\%}, and a median time-to-fix of \textbf{2.45\,s}. These results support the feasibility of real-time virtual mentoring for CS1, while larger controlled studies are still needed to evaluate comparative effectiveness and long-term learning impact.

Future work will focus on three directions: (i) comparative evaluation of the rule baseline against a learned policy selector trained from interaction logs, (ii) larger classroom-scale and longitudinal studies with student-level outcome measures, and (iii) richer learner-state modeling that preserves interpretability and bounded latency.

\section*{Acknowledgment}

This work was carried out as part of the master's degree program at Astana IT University and is connected to the development of the \system{} prototype for the author's thesis on virtual mentor models in programming education.

\bibliographystyle{IEEEtran}
\bibliography{references}

\end{document}
